\hypersetup{pageanchor=false}
\begin{titlepage}
\centering
{\large 일반국민/대학(원)생 부문}\par\vspace{4mm}
{\Large\bfseries 제6회 K-인공지능 제조데이터 분석 경진대회 보고서}\par\vspace{12mm}
\renewcommand{\arraystretch}{1.6}
\begin{tabular}{|p{0.18\linewidth}|p{0.72\linewidth}|}
\hline
\textbf{프로젝트명} & 생산계획 기반 조건부 베이지안 전달모형(C-BAT)을 이용한 제조공장 전력 예측과 최대피크 위험조건 분석 \\
\hline
\textbf{팀명} & \rule{5cm}{0.4pt} \\
\hline
\textbf{내용요약} &
\small 제조공장의 하루 96개 15분 슬롯 전력을 매일 00:00에 그날의 시간별 생산계획, 달력, 전날까지의 관측 전력만으로 예측하고, 최대피크 위험과 요금에 영향을 주는 날을 가려내는 방법을 제안한다.
제안 모형 C-BAT은 가동유형별 기본곡선, 과거 기준일 곡선, 생산계획의 전달항, 가동 여부에 따라 달라지는 잡음을 베이지안으로 함께 추정하고, 예측 경로에서 피크 분포와 초과확률을 계산한다.
완전·편집 복사일 122일을 목표에서 제외하고 시간순 검증을 설계하였다.
개발 구간(2021-07-07\,--\,08-31)과 9월 1--14일을 합친 평가는 67일·6,358슬롯이며 피크 평가일은 65일이다.
합산 피크 MAE는 C-BAT 8.93\,kW, 기존 BAT 13.09\,kW, LightGBM 기본 모형 M2 14.54\,kW로, C-BAT가 각각 31.8\%, 38.6\% 낮았다. 날짜 bootstrap 95\% 구간에서도 두 개선이 확인되었다.
합산 슬롯 MAE는 C-BAT 8.22\,kW, 기존 BAT 7.61\,kW, M2 10.57\,kW였다. 개발 구간 고피크일 13일의 피크 MAE는 C-BAT 4.37\,kW, 튜닝 LightGBM 18.27\,kW였다.
월 최대 갱신 경보는 개발 구간 백테스트에서 유일한 갱신일을 포착했으나 경보 10회 중 9회가 오경보였다. 금액 133,120원은 완벽 조치의 상한이며 실제 절감액은 아니다. \\
\hline
\end{tabular}
\vfill
\begin{flushleft}
\small 상기 본인(팀)은 위의 내용과 같이 제6회 K-인공지능 제조데이터 분석 경진대회 결과 보고서를 제출합니다.
\end{flushleft}
\vspace{6mm}
{2026년 \quad 10월 \quad\quad 일}\par\vspace{8mm}
\begin{flushright}
팀장 : \rule{3.5cm}{0.4pt} (서명)\\[3mm]
팀원 : \rule{3.5cm}{0.4pt} (서명)\\[3mm]
팀원 : \rule{3.5cm}{0.4pt} (서명)
\end{flushright}
\vspace{10mm}
{\large (사)중소기업기술혁신협회 귀중}
\end{titlepage}
\hypersetup{pageanchor=true}

\section*{평가 항목과 보고서 구성}
\addcontentsline{toc}{section}{평가 항목과 보고서 구성}
보고서의 장은 공식 평가 항목과 아래와 같이 대응한다.

\begin{center}
\small
\begin{tabular}{p{0.26\linewidth}p{0.25\linewidth}p{0.38\linewidth}}
\toprule
평가 항목 (배점) & 보고서 위치 & 주요 내용 \\
\midrule
1. 데이터 이해 및 진단 (15) & 3장 & 변수, 전처리 구조, 복사일·결측, 누수 진단, 검증 설계 \\
2. AI 예측모델 개발 및 성능평가 (40) & 2장, 4.1--4.2절, 4.6--4.7절 & 방법론, 비교 모델, 성능, 절제, 9월 시험, 최종 선택 \\
3. 영향요인 및 오류분석 (15) & 4.3--4.5절 & 오차·피크 발생 조건, 미탐지·오경보, 입력 오차 \\
4. 현장 활용방안 (10) & 5장 & 월 최대 갱신 경보, 백테스트, 현장 검증 \\
5. 창의성 및 차별성 (10) & 6.1절 & 선행 방법 대비 차별점과 효과 \\
6. 코드 및 재현성 (10) & 부록 A & 한 줄 실행, 재현 명령, 누수 점검 \\
\bottomrule
\end{tabular}
\end{center}
\clearpage
